
A neural network that achieves high predictive accuracy can fundamentally misunderstand the structure of its input, learning correct answers through an incorrect mechanism. This counterintuitive finding lies at the heart of understanding when architectural inductive biases are essential versus merely convenient.

Weight-space learning---predicting model properties directly from weight tensors---has emerged as a paradigm for tasks ranging from accuracy prediction to model quality assessment. Central to this paradigm is the question of how to handle permutation symmetries: the weights of a neural network can be permuted across hidden neurons without changing the function it computes. Permutation-equivariant architectures like Neural Functional Networks (NFN) explicitly encode this symmetry, while standard MLPs must learn it from data.

The conventional wisdom holds that with sufficient data diversity, non-equivariant architectures should eventually learn to exploit these symmetries. If the training distribution contains many permuted versions of similar weight configurations, an MLP should discover that permutation-related weights map to similar properties. This assumption has implications: if true, architectural inductive biases become a matter of convenience rather than necessity.

We challenge this assumption through systematic experiments across data scales from 1,000 to 40,000 models. Our finding: MLPs trained on 40,000 models achieve $R^{2}=0.986$ but fail to develop permutation invariance (invariance score 0.63 versus NFN's 1.0). This dissociates task performance from mechanism---high accuracy coexists with representations that exploit dataset-specific position-accuracy correlations rather than semantic weight structure.

Our contributions are:

\begin{enumerate}
\item \textbf{Sample Efficiency of Equivariance:} NFN achieves $R^{2}=0.95$ at N=1,000 while matched-capacity MLPs achieve $R^{2}=0.35$---a 60 percentage point advantage.
\end{enumerate}

\begin{enumerate}
\item \textbf{Mechanism Verification:} NFN's invariance is near-perfect (deviation < 1.$19\times 10^{-7}$ across permutations) while untrained MLPs show coefficient of variation 0.19 under the same permutations.
\end{enumerate}

\begin{enumerate}
\item \textbf{Falsification of Learned Invariance:} MLPs cannot learn permutation invariance from data diversity. At N=40,000, MLPs achieve $R^{2}=0.986$ but invariance=0.63, below the 0.8 threshold.
\end{enumerate}

